%% GENERATED by python/paper/build.py from results/ -- do not edit by hand.
%% Rebuild:  .venv\Scripts\python.exe python\paper\build.py --only UK_ESC_FrenchAll
%% Source:   results/esc/escExperimentsUK.csv
\begin{table}[!htb]
\centering
\begin{threeparttable}
\caption{Endogenous design and all French characteristics in the UK}
\label{table:UK_ESC:frenchAll}
\renewcommand{\arraystretch}{1.25}
\begin{tabularx}{\textwidth}{p{2.6cm}YYYYY}
\toprule
\textbf{Scenario} & \textbf{CRRA} ($\rho$) & $\bm{\theta}$ \textbf{(2020)} & \textbf{Tax rate} & \textbf{Savings rate} & \textbf{Avg. workweek} \\
\midrule 
 & 0.5 & 0.560 & 11.86\% & 19.11\% & 34.73 \\ % row: baseline@0.5
Baseline & 1.0 & 0.560 & 11.86\% & 17.25\% & 34.75 \\ % row: baseline@1.0
 & 2.0 & 0.560 & 11.86\% & 15.81\% & 34.74 \\[.5em]\hline\\[-.75em] % row: baseline@2.0
 & 0.5 & 0.560 & 11.97\% & 0.00 p.p. & 35.75 \\ % row: pinned@0.5
Exogenous $\theta$ & 1.0 & 0.560 & 12.31\% & $-0.14$ p.p. & 35.66 \\ % row: pinned@1.0
 & 2.0 & 0.560 & 12.63\% & $-0.17$ p.p. & 35.63 \\[.5em]\hline\\[-.75em] % row: pinned@2.0
 & 0.5 & 0.564 & 11.98\% & $-0.05$ p.p. & 35.76 \\ % row: chosen@0.5
Endogenous $\theta$ & 1.0 & 0.521 & 12.30\% & $-0.13$ p.p. & 35.65 \\ % row: chosen@1.0
 & 2.0 & 0.414 & 12.64\% & $-0.12$ p.p. & 35.53 \\[.5em]\hline\\[-.75em] % row: chosen@2.0
 & 0.5 & 1.000 & 21.29\% & $-4.45$ p.p. & 35.24 \\ % row: france@0.5
France & 1.0 & 1.000 & 21.29\% & $-3.81$ p.p. & 35.24 \\ % row: france@1.0
 & 2.0 & 1.000 & 21.29\% & $-3.34$ p.p. & 35.24 \\[.5em]\hline\\[-.75em] % row: france@2.0
\bottomrule
\end{tabularx}
\begin{tablenotes}[flushleft]
\footnotesize
\item[] \textit{Note:} The cost specification, the cost parameter and the units are those of table \ref{table:UK_ESC:incomeDistr}. The scenario replaces the UK $\eta_i$, the level of $X_i$ and the voting weights $\mu_i$ with France\textquotesingle s simultaneously. The France row is France's own calibrated path rather than a counterfactual on the UK model: it carries France\textquotesingle s own $\omega$ as well as its characteristics, and its workweek is a calibration target rather than a prediction.
\end{tablenotes}
\end{threeparttable}
\end{table}
